(1) 北半球纬度 \(\lambda\) 的地面实验室有长度 \(L\) 的小振幅Foucault摆，局地有效重力加速度 \(g\)，地球自转角速度 \(\Omega_\oplus\)。以东向 \(x\)、北向 \(y\)、竖直向上 \(z\) 为右手坐标，令 \(\omega_0=\sqrt{g/L}\)。在 \(\Omega_\oplus\ll\omega_0\) 下忽略水平地转分量造成的高阶竖直耦合及空气阻尼。写出水平线性运动方程与复坐标方程。
(2) 求两圆偏振本征频率，说明其劈裂与地球自转垂直分量的关系。
(3) 对初始近似平面振动，求振动平面相对地面的慢进动角速度与完整进动周期，分别检查赤道和北极极限。
(4) 由运动方程直接证明实验室内摆的小振幅机械能守恒，解释科里奥利力虽改变振动方向却不做功，并讨论近似的适用时间尺度。
